% Generated by scripts/build_blueprint_tex.py from data/model.json -- do not edit by hand.
\begin{bpstratalegend}
\bpstratumrow{arith}{Galois representations, theta lifts and Iwasawa-type cohomology bounds that are meant to produce the families on which a prime sum is small.}{\bpcompref{cohomology}, \bpcompref{primes}, \bpcompref{galois}, \bpcompref{inputs}}
\bpstratumrow{est}{A finite character comparison turns geometric control of a chain of theta lifts of length about \(\sqrt{x}\) into an upper bound for a weighted prime sum.}{\bpcompref{comparison}, \bpcompref{family}, \bpcompref{hosts}}
\bpstratumrow{ana}{A localized Landau oscillation argument converts one-sided bounds on whole power windows into a zero-free half-plane.}{\bpcompref{statistic}, \bpcompref{residues}, \bpcompref{criterion}, \bpcompref{choreography}}
\bpstratumrow{dest}{Euler factors, conductors and functional equations turn the half-plane theorem into the five exported statements, including Mathlib's \href{https://github.com/leanprover-community/mathlib4/blob/0df444a360eaa60ab8c11dca51a86af692955474/Mathlib/NumberTheory/LSeries/RiemannZeta.lean\#L182-L186}{\bpcode{RiemannHypothesis}}.}{\bpcompref{endpoints}, \bpcompref{reflection}, \bpcompref{descent}}
\end{bpstratalegend}
